\documentclass[conference]{IEEEtran}

\usepackage[utf8]{inputenc}
\usepackage{amsmath,amssymb,amsfonts,bm}
\usepackage{booktabs}
\usepackage{multirow}
\usepackage{array}
\usepackage{enumitem}
\usepackage{algorithm}
\usepackage{algpseudocode}
\usepackage{graphicx}
\usepackage{xcolor}
\usepackage{url}
\usepackage[hidelinks]{hyperref}

\DeclareMathOperator*{\argmin}{arg\,min}
\DeclareMathOperator*{\argmax}{arg\,max}
\DeclareMathOperator{\softmin}{softmin}
\DeclareMathOperator{\softmaxmax}{softmaxmax}
\DeclareMathOperator{\sigmoid}{sigmoid}
\DeclareMathOperator{\dist}{dist}
\DeclareMathOperator{\Quantile}{Quantile}
\DeclareMathOperator{\softmax}{softmax}
\DeclareMathOperator{\WQ}{WQ}
\DeclareMathOperator{\UCB}{UCB}
\DeclareMathOperator{\LCVaR}{LCVaR}

\newcommand{\E}{\mathbb{E}}
\newcommand{\cA}{\mathcal{A}}
\newcommand{\cB}{\mathcal{B}}
\newcommand{\cC}{\mathcal{C}}
\newcommand{\cD}{\mathcal{D}}
\newcommand{\cF}{\mathcal{F}}
\newcommand{\cG}{\mathcal{G}}
\newcommand{\cI}{\mathcal{I}}
\newcommand{\cJ}{\mathcal{J}}
\newcommand{\cL}{\mathcal{L}}
\newcommand{\cM}{\mathcal{M}}
\newcommand{\cO}{\mathcal{O}}
\newcommand{\cP}{\mathcal{P}}
\newcommand{\cR}{\mathcal{R}}
\newcommand{\cS}{\mathcal{S}}
\newcommand{\cT}{\mathcal{T}}
\newcommand{\cU}{\mathcal{U}}
\newcommand{\cZ}{\mathcal{Z}}
\newcommand{\calX}{\mathcal{X}}
\newcommand{\rvn}{\mathrm{OC\mbox{-}MERO}}
\newcommand{\cpdm}{\mathrm{ReCoT}}
\newcommand{\rec}{\mathrm{rec}}
\newcommand{\drv}{\mathrm{drv}}
\newcommand{\exec}{\mathrm{exec}}
\newcommand{\near}{\mathrm{near}}
\newcommand{\contact}{\mathrm{contact}}
\newcommand{\imp}{\mathrm{imp}}
\newcommand{\terminal}{\mathrm{terminal}}
\newcommand{\route}{\mathrm{route}}
\newcommand{\wit}{\mathrm{wit}}
\newcommand{\harm}{\mathrm{harm}}
\newtheorem{proposition}{Proposition}

\title{Observation-Consistent Recoverability as a Calibrated Planning Primitive for Autonomous Driving}

\author{Anonymous Authors}

\begin{document}
\maketitle

\begin{abstract}
Autonomous-driving planners are often optimized for imitation, progress, comfort, or immediate collision risk.  Such objectives can appear safe locally while consuming the space, control authority, or scene topology needed to recover from later uncertainty or post-contact dynamics.  We propose OC-RAP, an observation-consistent recoverability planner that evaluates short executable prefixes by the recovery headroom they preserve.  OC-RAP learns root-shared counterfactual evidence with ReCoT, maps signed recovery margins through a monotone, observation-consistent existential operator (OC-MERO), and selects actions with CRISP, a calibrated recovery-preserving selector.  The key distinction is that recovery witnesses must be selectable from post-prefix observations, avoiding oracle labels that choose a different hidden witness for each indistinguishable latent future.  We specify recovery labels, calibration constraints, and evaluation metrics for normal, low-headroom, near-contact, and post-contact driving, and provide closed-loop and offline experiment templates for recovery success, false admission, harm non-inferiority, minimal intervention, robustness, and runtime.
\end{abstract}

\section{Introduction}
Autonomous driving safety is often framed as avoiding the next collision or minimizing a risk-weighted trajectory cost.  This framing is incomplete.  A locally safe action can still erase future braking authority, lateral escape space, route topology, or post-contact stabilization headroom.  Conversely, two actions with similar immediate collision exposure may differ substantially in whether the ego vehicle can later stop safely, rejoin the route, yield into a gap, avoid a secondary collision, or stabilize after unavoidable contact.  Recent analyses of robotaxi failures and post-impact vehicle-control studies underline that the consequences of an event are shaped not only by whether contact occurs, but also by the control authority and environment structure left available after the preceding maneuver~\cite{koopman2024anatomy,yang2025post,ghosh2026post}.

Existing planning and safety tools address related pieces of this problem, but not the full normal-to-critical continuum.  Reachability and viability methods provide formal semantics for safe or reach-avoid sets~\cite{bansal2017hamilton,bansal2021deepreach,lin2023generating}; control barrier functions and predictive safety filters enforce invariance or backup feasibility for proposed controls~\cite{ames2019control,koller2018learning,wabersich2021predictive,tearle2021predictive,kim2026backup}; risk-aware and contingency planners reason over uncertain futures with chance constraints, CVaR, or scenario trees~\cite{uryasev2000conditional,hakobyan2021wasserstein,safaoui2024distributionally,li2023marc,mustafa2024racp}; and post-impact methods focus on stability and secondary-collision mitigation once an event is imminent or has occurred~\cite{bosia2024real,wang2023integrated,parseh2023motion,wang2024post}.  These approaches motivate recovery reasoning, but they usually treat recoverability as either a terminal backup check, a scalar penalty, or a regime-specific emergency module.

We argue that recoverability should be a first-class action property.  The relevant question is counterfactual: after executing a short candidate prefix, does the resulting scene still admit an executable recovery witness under plausible future roots?  This question is existential over recovery witnesses but distributional over latent futures.  A single good recovery option is sufficient for a particular future only if the option can be selected from information actually observable after the prefix.  Otherwise, an oracle label may silently choose different hidden witnesses for latent futures that induce the same post-prefix observation, overstating deployable recoverability.

OC-RAP addresses this issue by combining learned evidence with structured recovery semantics.  Recovery specifications are represented as signed margins over no-contact viable recovery, minimum-risk contact-imminent recovery, and post-contact stabilization; the reported option margin is the largest specification margin available to the witness, not the minimum across mutually exclusive regimes.  Witness selection is constrained by post-prefix observation equivalence classes, so indistinguishable latent roots share a witness distribution.  The resulting lower-tail recoverability score can be calibrated as an admissibility constraint rather than tuned as another cost weight.

This paper makes three contributions:
\begin{enumerate}[leftmargin=*]
    \item \textbf{Observation-consistent recoverability.} We formulate recoverability as a deployable, post-prefix property: recovery witnesses are existential, but they must be chosen consistently over latent futures that are observationally indistinguishable after the prefix.
    \item \textbf{Structured recovery inference.} We introduce ReCoT and OC-MERO, which predict signed recovery margins under root-shared counterfactual futures and aggregate them with a monotone, $\mu$-weighted existential lower-tail operator instead of a black-box scalar value head.
    \item \textbf{Calibrated recovery-preserving selection.} We introduce CRISP, a selector that preserves nominal utility when it is already admissible, and otherwise switches using calibrated constraints on recoverability, absolute harm, relative harm non-inferiority, and rule/hard-shell violation.
\end{enumerate}

\section{Related Work}

\subsection{Collision Mitigation, Unavoidable Collision Planning, and Post-Impact Control}
Collision mitigation and post-impact control study how autonomous or assisted vehicles should behave when collision risk is high or impact has already occurred. Prior work has optimized collision severity, impact configuration, and post-impact motion, including unavoidable collision mitigation~\cite{li2022vehicle,bosia2024real}, post-impact planning~\cite{wang2023integrated,parseh2023motion}, and post-impact stability or trajectory restoration control~\cite{yang2025post,wang2024post,ghosh2026post}. These methods motivate the importance of post-event controllability and secondary-collision avoidance. However, they usually focus on emergency or post-impact regimes. In contrast, we treat post-event recovery as one regime of a broader recoverability principle. Our planner evaluates each candidate action by the recovery headroom it preserves before the event occurs, while using first-impact harm non-inferiority to avoid trading immediate harm for easier recovery.

\subsection{Viability, Reachability, CBFs, Predictive Safety Filters, and Backup-Based Safety Filters}
Viability and reachability provide formal tools for characterizing states from which safety can be maintained under admissible controls~\cite{bansal2017hamilton,bansal2021deepreach,lin2023generating}. CBFs enforce forward invariance of safe sets through optimization-based controllers~\cite{ames2019control}. Predictive safety filters and backup-based filters evaluate whether proposed controls can be safely applied, often relying on backup policies, terminal sets, or model-predictive shielding~\cite{wabersich2021predictive,tearle2021predictive,koller2018learning,kim2026backup}. These works provide the conceptual foundation for return-to-safe-set reasoning. Our method differs in scope and mechanism: rather than online filtering of a single proposed control, we learn an action-conditioned, distributional recoverability margin over many candidate actions in high-dimensional interactive driving scenes. The learned margin is used as a lower-confidence gate and ranking signal while preserving nominal utility whenever recovery headroom is sufficient.

\subsection{Risk-Aware, Contingency, and World-Model-Based Planning}
Risk-aware planning optimizes under uncertainty using chance constraints, conditional value-at-risk (CVaR), distributionally robust optimization, or other tail-risk objectives~\cite{uryasev2000conditional,hakobyan2021wasserstein,safaoui2024distributionally,nair2024predictive}. Contingency planners construct policy-conditioned or scenario-tree futures to handle multi-modal interactions~\cite{li2023marc,mustafa2024racp}. World-model and transformer-based planners learn compact representations for predicting future agents, interactions, or planning costs~\cite{vaswani2017attention,ding2025understanding,hu2025survey}. Our ReCoT follows this lineage but changes the prediction target: it predicts recovery-sufficient evidence rather than visually faithful or trajectory-accurate futures. OC-MERO then maps this evidence to observation-consistent, closed-loop teacher-supervised return-to-viable-set margins rather than generic collision risk.





% =========================
% Suggested title
% =========================
% \title{Observation-Consistent Recoverability as a Calibrated Planning Primitive for Autonomous Driving}

\section{Method}
\label{sec:method}

The planner evaluates a short executable ego prefix by the recovery headroom it leaves, not by the likelihood or cost of the prefix alone.  Given belief $b_t$, route command $q_t$, and a prefix $a$, the central question is whether the post-prefix scene admits an executable recovery witness that can be selected from information available after executing $a$.  The method has four stages: candidate prefixes are paired with scene-conditioned recovery affordance tokens; a recoverability-centered transformer predicts root-shared counterfactual evidence; an observation-consistent witness policy converts evidence into a monotone existential lower-tail recoverability score; and a calibrated selector chooses the least intrusive admissible action.

\subsection{Prefix Recoverability and Executable Recovery Affordances}
\label{subsec:prefix_affordances}

The input belief is
\begin{equation}
\label{eq:method_belief_new}
    b_t=(\mathcal I_{t-H_h:t}, q_t, \Sigma_t),
\end{equation}
where $\mathcal I$ contains ego-centered map, route, traffic-control, agent-history, occupancy, occlusion, and uncertainty channels.  A proposal module produces prefix parameters $\xi_k$ and projects them to dynamically executable prefixes
\begin{equation}
\label{eq:method_action_set_new}
    \mathcal A_0(b_t)=\{a_1,\ldots,a_K\},\qquad a_k=\Gamma(\xi_k).
\end{equation}
The projection $\Gamma$ enforces ego dynamics, input limits, static swept-geometry constraints, and route-compatible heading bounds.  A prefix is intentionally short; it is the action to execute before the next replanning step, not a complete future plan.

Recoverability is represented by executable recovery affordance tokens rather than a scalar critic.  For prefix $a$, the token set is
\begin{equation}
\label{eq:method_affordance_set}
    \mathcal W_a(b_t)=\{w_j\}_{j=1}^{L_a},\qquad
    w_j=(\rho_j,d_j,\kappa_j,\psi_j,\mathcal H_j,\chi_j).
\end{equation}
Here $\rho_j$ is a scene anchor such as a lane tube, legal stop region, yield boundary, free-space pocket, route-rejoin target, or stabilization basin; $d_j$ is an interpretable affordance tag used only for analysis and stratification; $\kappa_j$ is the fallback controller attached to the token; $\psi_j$ parameterizes a local recovery potential; $\mathcal H_j$ is a hard-shell constraint set; and $\chi_j$ stores token metadata such as reachable horizon, controller class, and anchor geometry.  The token is therefore a candidate recovery witness: it specifies where to recover, what constraints must hold, and which fallback controller can be executed.

A token is admissible only through its hard shell and conservative recovery potential.  The hard shell contains physical feasibility, rule compliance outside minimum-risk fallback, and stability margins.  The learned residual can change the local recovery potential but cannot override hard-shell violations.  The teacher evaluates three mutually exclusive recovery specifications: no-contact viable recovery, minimum-risk contact-imminent recovery, and post-contact stabilization.  The option-level margin is the best specification available to the same witness,
\begin{equation}
\label{eq:method_margin}
    G_{a,j}^{m,\star}=\max_{s\in\{\mathrm{no},\mathrm{mr},\mathrm{post}\}}
    G_{s,a,j}^{m,\star},
\end{equation}
not the minimum across all specifications.  This avoids penalizing a no-contact witness by an irrelevant post-contact condition, or a post-contact stabilization witness by an irrelevant no-contact condition.  The selected specification identity is stored for supervision and diagnostics, but it is not used as an online regime switch.  Appendix~\ref{app:recovery_affordances} gives the exact hard-shell margins, residual bound, fallback controllers, and token generator.

\subsection{ReCoT: Root-Shared Counterfactual Evidence Prediction}
\label{subsec:recot}

The network predicts recovery evidence rather than decoding a full world model.  For each root scene, root uncertainty seeds $\omega^m$ are sampled once and shared across all candidate prefixes:
\begin{equation}
\label{eq:method_root_futures}
    \omega^m\sim p_\theta(\omega\mid b_t),\qquad
    z_a^m=T_\theta(b_t,a,\omega^m),\qquad
    \hat p^m=p_\theta(m\mid b_t).
\end{equation}
Root sharing makes the comparison between two prefixes counterfactual: action differences are evaluated under the same latent interaction roots, simulator perturbations, and occlusion releases.

The recoverability-centered transformer, ReCoT, uses four token families: scene tokens from vector and BEV encoders, prefix tokens from sampled ego poses and controls, root tokens from $\omega^m$, and recovery-affordance tokens from $w_j$.  Cross-attention is performed over the tensor product of prefix, root, and witness tokens, yielding
\begin{equation}
\label{eq:method_recot_tuple}
    h_{a,j}^{m}=\mathrm{ReCoT}(h_b,h_a,h_{\omega^m},h_{w_j},\ell_{a,j}),
\end{equation}
where $\ell_{a,j}$ pools swept-corridor and anchor-overlap features.  The evidence head outputs
\begin{equation}
\label{eq:method_evidence_new}
    \widehat e_{a,j}^{m}=\left(
    \widehat{\bm g}_{a,j}^{m},
    \widehat y_{a,j}^{m},
    \widehat h_{a}^{m},
    \widehat k_{a,j}^{m},
    \widehat u_{a,j}^{m},
    \widehat c_{a}^{m},
    \widehat\beta_a^m
    \right).
\end{equation}
The vector $\widehat{\bm g}$ contains signed recovery-specification margins; $\widehat y$ is witness feasibility; $\widehat h$ is first-contact harm exposure; $\widehat k$ is post-contact recovery deficit; $\widehat u$ is epistemic or root-tail uncertainty; $\widehat c$ is the non-slacked rule or hard-shell violation used by CRISP; and $\widehat\beta$ is the predicted post-prefix posterior.  The evidence loss is
\begin{equation}
\label{eq:method_evidence_loss}
\begin{split}
\mathcal L_{\mathrm{ev}}=
&\lambda_g\ell_{\mathrm{Huber}}(\widehat{\bm g},\bm g^\star)
+\lambda_y\ell_{\mathrm{BCE}}(\widehat y,Y^\star)
+\lambda_h\ell_{\mathrm{pin}}(\widehat h,H^\star)\\
&+\lambda_k\ell_{\mathrm{pin}}(\widehat k,K^\star)
+\lambda_u\ell_{\mathrm{NLL}}(\widehat u,U^\star)
+\lambda_c\ell_{\mathrm{Huber}}(\widehat c,C^\star)
+\lambda_b D_{\mathrm{KL}}(\beta^\star\Vert\widehat\beta).
\end{split}
\end{equation}
All terms are computed on teacher rollouts over retained $(a,j,m)$ tuples.  The teacher construction is specified in Appendix~\ref{app:teacher_labels_new}.

\subsection{OC-MERO: Observation-Consistent Existential Recovery}
\label{subsec:oc_mero}

A purely existential recovery label can leak oracle information: the planner may claim that a different witness exists in each latent future even when those futures are indistinguishable after the prefix.  We therefore restrict witness selection to post-prefix information.  Executing $a$ under root future $m$ produces observation $o_a^m$ and posterior
\begin{equation}
\label{eq:method_posterior_new}
    \beta_a^m=q_\theta(z\mid b_t,a,o_a^m).
\end{equation}
The witness policy is
\begin{equation}
\label{eq:method_witness_policy}
    \mu_\omega(w_j\mid \widehat\beta_a^m,a)=
    \mathrm{softmax}_j\,q_\omega(w_j,\widehat\beta_a^m,a).
\end{equation}
If two roots induce equivalent post-prefix observations, the policy must assign the same witness distribution.  This is enforced by
\begin{equation}
\label{eq:method_obs_loss}
    \mathcal L_{\mathrm{obs}}=
    \sum_{a,m,m'}\mathbf 1[d_{\mathrm{obs}}(o_a^m,o_a^{m'})\le\epsilon_o]
    D_{\mathrm{JS}}(\mu_a^m\Vert\mu_a^{m'}).
\end{equation}

Each witness is scored by a monotone map
\begin{equation}
\label{eq:method_option_score_new}
    v_{a,j}^{m}=f_{\phi}^{\mathrm{mono}}(
    \widehat{\bm g}_{a,j}^{m},
    -\widehat h_a^m,
    -\widehat k_{a,j}^m,
    -\widehat u_{a,j}^m),
\end{equation}
so improving any safety margin cannot reduce recoverability, and increasing harm, post-contact deficit, or uncertainty cannot improve it.  Observation-consistent existential recovery for mode $m$ is
\begin{equation}
\label{eq:method_existential_new}
    \pi_a^m=\sigmoid\!\left(
    \tau_R\log\sum_{j\in\mathcal V_a}
    \mu_\omega(w_j\mid\widehat\beta_a^m,a)
    \exp(v_{a,j}^{m}/\tau_R)-c_R
    \right),
\end{equation}
where $\mathcal V_a$ is the valid token set.  Since $\mu_\omega(\cdot\mid\widehat\beta_a^m,a)$ is normalized over valid tokens, the default offset is $c_R=0$; using an additional $\tau_R\log L_a$ normalization would double-count option-count normalization and artificially depress prefixes with more valid witnesses.  The operator approaches the best implementable witness as $\tau_R\to0$.  The action-level recoverability used by the selector is the lower-tail value
\begin{equation}
\label{eq:method_tail_recovery_new}
    R_a=\mathrm{LCVaR}_{\alpha_R}(\{\pi_a^m,\hat p^m\}_{m=1}^{M}).
\end{equation}
OC-MERO also reports bottleneck pressure $B_a$, tail uncertainty $U_a$, harm $H_a$, relative harm gap $\Delta H_a$, post-contact deficit $K_a$, non-slacked rule or hard-shell violation $C_a$, selected witness $\hat w_a$, and witness concentration $W_a$:
\begin{equation}
\label{eq:method_profile_new}
    \Pi_a=(R_a,B_a,U_a,H_a,\Delta H_a,K_a,C_a,\hat w_a,W_a).
\end{equation}
The operator loss combines teacher lower-tail recoverability, same-root ranking, observation-consistent witness supervision, and sampled monotonicity checks:
\begin{equation}
\label{eq:method_operator_loss}
    \mathcal L_{\mathrm{op}}=
    \lambda_R\ell_{\mathrm{pin}}(R_a,R_a^\star)
    +\lambda_{\mathrm{rank}}\ell_{\mathrm{rank}}
    +\lambda_w\ell_{\mathrm{CE}}(\hat w_a,w_a^\star)
    +\lambda_o\mathcal L_{\mathrm{obs}}
    +\lambda_m\mathcal L_{\mathrm{mono}}.
\end{equation}

\subsection{CRISP: Calibrated Recovery-Preserving Selection}
\label{subsec:crisp_new}

The selector treats recoverability as a calibrated admissibility condition, not as a tunable cost term.  Candidate prefixes may be refined for a small number of projected gradient or SQP steps,
\begin{equation}
\label{eq:method_refinement_new}
    \max_{\xi}\; U_{\mathrm{drv}}(b_t,\Gamma(\xi))
    -\lambda_s\|\xi-\xi_0\|_2^2
    \quad\mathrm{s.t.}\quad \Gamma(\xi)\in\mathcal H_{\mathrm{dyn}},
\end{equation}
then re-evaluated by ReCoT and OC-MERO.  Calibration offsets $q=(q_R,q_H,q_\Delta,q_C)$ are computed on a held-out calibration split using conformal risk control and conformal decision theory~\cite{angelopoulos2024conformal,lekeufack2024conformal}.  The calibrated admissible set is
\begin{equation}
\label{eq:method_admissible_set}
\begin{split}
    \mathcal S(q)=\{a\in\mathcal A:
    & R_a-q_R\ge\eta_R,
      H_a+q_H\le\eta_H,\\
    & \Delta H_a+q_\Delta\le\epsilon_H,
      C_a+q_C\le0\},
\end{split}
\end{equation}
where $C_a$ is the predicted non-slacked rule violation outside the minimum-risk fallback.

Let $a_{\mathrm{nom}}=\argmax_{a\in\mathcal A}U_{\mathrm{drv}}(b_t,a)$.  If $a_{\mathrm{nom}}\in\mathcal S(q)$, the planner executes $a^\star=a_{\mathrm{nom}}$.  If the nominal prefix is not admissible and $\mathcal S(q)$ is nonempty,
\begin{equation}
\label{eq:method_constrained_selector}
    a^\star=\argmax_{a\in\mathcal S(q)}
    \left[U_{\mathrm{drv}}(b_t,a)-\lambda_BB_a-\lambda_UU_a-\lambda_KK_a\right].
\end{equation}
If $\mathcal S(q)=\emptyset$, the planner returns a minimum-risk prefix
\begin{equation}
\label{eq:method_relaxation}
    a^\star=\argmin_{a\in\mathcal A}
    \left[
    \lambda_R[\eta_R-(R_a-q_R)]_+
    +\lambda_H[H_a+q_H-\eta_H]_+
    +\lambda_\Delta[\Delta H_a+q_\Delta-\epsilon_H]_+
    +\lambda_KK_a+
    \lambda_C[C_a+q_C]_+
    \right].
\end{equation}
The executed fallback controller is the controller $\kappa_{\hat w_a}$ attached to the selected witness when a witness is available, and otherwise the minimum-risk stop/stabilize controller.  The proposal and utility heads use
\begin{equation}
\label{eq:method_selector_loss_new}
    \mathcal L_{\mathrm{sel}}=
    \lambda_{\mathrm{imit}}\ell_{\mathrm{NLL}}(a_{\mathrm{exp}}\mid b_t)
    +\lambda_{\mathrm{util}}\ell_{\mathrm{rank}}(U_{\mathrm{drv}},U_{\mathrm{teach}})
    +\lambda_{\mathrm{proj}}\ell_{\mathrm{dyn}}.
\end{equation}
The total training objective is
\begin{equation}
\label{eq:method_total_loss_new}
    \mathcal L_{\mathrm{total}}=\mathcal L_{\mathrm{ev}}+\mathcal L_{\mathrm{op}}+\mathcal L_{\mathrm{sel}}.
\end{equation}
Calibration is performed after training and does not update network weights.


\section{Experiments}
\label{sec:experiments}

The experiments are designed to test whether observation-consistent recoverability improves closed-loop behavior, whether the learned recovery profiles match teacher labels, and whether calibration controls selected-action false admission under both in-distribution and shifted conditions.  All numerical entries in the tables are left as placeholders so that results can be filled after running the benchmark.

\subsection{Benchmark and Evaluation Protocol}
\label{subsec:benchmark_protocol}

We evaluate on \textsc{MetaDrive-Recovery}, a counterfactual closed-loop benchmark built on MetaDrive~\cite{li2022metadrive} with optional WOMD replay and hybrid stress roots from ScenarioNet-style conversion~\cite{ettinger2021large,li2023scenarionet}.  A root scene stores the BEV belief $b_i$, route command $q_i$, candidate prefixes $\mathcal A_0(b_i)$, root-shared latent futures $\{\omega_i^m\}_{m=1}^M$, recovery-affordance tokens $\mathcal W_a(b_i)$, and teacher labels for each retained $(a,j,m)$ tuple.  Root seeds are shared across all candidate prefixes so that same-root action ranking compares actions under paired future uncertainty rather than independent future draws.

Root scenes are balanced across normal, low-headroom, near-contact, and contact/post-contact tags.  These tags are used only for balancing, analysis, and Mondrian calibration; the online planner does not use a regime gate.  In closed loop, all planners receive the same belief history, route command, map, actor states, and perturbation seeds.  The simulator executes the first segment of the selected prefix and replans.  Crash termination is disabled for contact/post-contact tests so that post-contact stabilization, secondary collision, safe stop, and route rejoin can be measured.

\paragraph{Observation-consistent labels.}
For each root $i$, prefix $a$, token $j$, and latent mode $m$, the teacher computes specification margins
\begin{equation}
\label{eq:exp_spec_margin}
    \bm G_{i,a,j}^{m,\star}=
    (G_{\mathrm{no}},G_{\mathrm{mr}},G_{\mathrm{post}})_{i,a,j}^{m,\star}.
\end{equation}
The option margin and option label are
\begin{align}
\label{eq:exp_option_margin_oc}
    G_{i,a,j}^{m,\star}
    &=\max_{s\in\{\mathrm{no},\mathrm{mr},\mathrm{post}\}}G_{s,i,a,j}^{m,\star},\\
\label{eq:exp_option_label_oc}
    Y_{i,a,j}^{m,\star}&=\mathbf 1[G_{i,a,j}^{m,\star}\ge0].
\end{align}
Let $c_{i,a}(m)$ denote the post-prefix observation-equivalence class of mode $m$.  The class-consistent teacher witness is
\begin{equation}
\label{eq:exp_class_witness_oc}
    j_{i,a,c}^{\star}=
    \argmax_{j\in\mathcal W_a(b_i)}
    \sum_{m:c_{i,a}(m)=c}\hat p_i^m Y_{i,a,j}^{m,\star},
\end{equation}
with ties broken by larger mean margin, smaller harm, smaller post-contact deficit, and smaller controller effort.  The deployable teacher mode label and lower-tail recoverability are
\begin{align}
\label{eq:exp_oc_mode_label}
    Y_{i,a,\mathrm{oc}}^{m,\star}
    &=Y_{i,a,j_{i,a,c_{i,a}(m)}^\star}^{m,\star},\\
\label{eq:exp_true_recoverability}
    R_{i,a}^{\star}
    &=\LCVaR_{\alpha_R}
    \left(\{Y_{i,a,\mathrm{oc}}^{m,\star},\hat p_i^m\}_{m=1}^{M}\right).
\end{align}
For diagnostics only, we also report the non-deployable oracle recoverability
\begin{equation}
\label{eq:exp_oracle_recoverability}
    R_{i,a}^{\mathrm{oracle}}=
    \LCVaR_{\alpha_R}
    \left(\{\max_j Y_{i,a,j}^{m,\star},\hat p_i^m\}_{m=1}^{M}\right),
\end{equation}
which can overestimate recoverability by selecting a different hidden witness for indistinguishable modes.

\begin{table}[t]
\centering
\caption{MetaDrive-Recovery root-scene split.  Splits are by root scene and original scenario identifier, not by $(b,a,j,m)$ tuple.}
\label{tab:dataset_split}
\begin{tabular}{lccccc}
\toprule
Split & Roots & Normal & Low-headroom & Near-contact & Contact \\
\midrule
Train & 20000 & 5000 & 5000 & 5000 & 5000 \\
Validation & 2000 & 500 & 500 & 500 & 500 \\
Calibration & 4000 & 1000 & 1000 & 1000 & 1000 \\
Test-ID & 4000 & 1000 & 1000 & 1000 & 1000 \\
Test-Shift & 4000 & 1000 & 1000 & 1000 & 1000 \\
\bottomrule
\end{tabular}
\end{table}

\begin{table}[t]
\centering
\caption{Default evaluation parameters.}
\label{tab:dataset_params}
\begin{tabular}{lcc}
\toprule
Parameter & Value & Meaning \\
\midrule
BEV extent & $80\,\mathrm{m}\times80\,\mathrm{m}$ & Ego-centered crop \\
BEV resolution & $256\times256$ & Main raster resolution \\
History horizon $H_h$ & $5$ & $1.0\,\mathrm{s}$ history \\
Prefix horizon $H_p$ & $10$ & $2.0\,\mathrm{s}$ prefix \\
Recovery horizon $H_r$ & $25$ & $5.0\,\mathrm{s}$ recovery \\
Candidate prefixes $K$ & $32$ & After projection/pruning \\
Latent modes $M$ & $8$ & Root-shared future seeds \\
Recovery tokens $L$ & $\le12$ & Valid tokens per prefix \\
Lower-tail mass $\alpha_R$ & $0.20$ & Recoverability LCVaR \\
\bottomrule
\end{tabular}
\end{table}

\subsection{Baselines}
\label{subsec:baselines}

Table~\ref{tab:baseline_definitions} lists all comparison methods.  All learned baselines use the same candidate prefix set unless stated otherwise; this isolates the effect of recovery inference and selection rather than proposal coverage.

\begin{table*}[t]
\centering
\footnotesize
\caption{Baseline planners and abbreviations.  The oracle selector is diagnostic and not deployable.}
\label{tab:baseline_definitions}
\begin{tabular}{p{0.12\linewidth}p{0.24\linewidth}p{0.54\linewidth}}
\toprule
Abbrev. & Method & Description \\
\midrule
NOM & Nominal utility planner & Selects the highest imitation/progress/comfort utility prefix without recovery constraints. \\
Risk-CVaR & Risk-aware CVaR planner~\cite{uryasev2000conditional,hakobyan2021wasserstein,safaoui2024distributionally} & Penalizes collision probability, harm, and tail risk with a scalar cost. \\
PSF-Backup & Predictive/backup safety filter~\cite{koller2018learning,wabersich2021predictive,tearle2021predictive,kim2026backup} & Accepts the nominal prefix only if a fixed brake-to-stop or lane-keep backup reaches a terminal safe set. \\
Contingency & Contingency planner~\cite{li2023marc,mustafa2024racp} & Scores policy-conditioned futures or scenario-tree branches and selects the best contingency-aware prefix. \\
WM-Recovery & World-model recovery scorer~\cite{ding2025understanding,hu2025survey} & Decodes future trajectory/occupancy samples and computes recovery features from decoded futures. \\
Scalar-R & Direct scalar critic & Predicts $R_a$ directly from $(b_t,a)$ without signed evidence or OC-MERO aggregation. \\
Viable-Hand & Hand-crafted viable-set critic~\cite{bansal2017hamilton,ames2019control} & Uses analytic distance-to-stop, lane-center, road-boundary, and free-space margins instead of learned evidence. \\
Oracle-OC & Observation-consistent oracle & Selects using teacher $R_{i,a}^{\star}$ and $H,C,K$ labels.  This upper bound is not deployable. \\
Ours & OC-RAP & Full ReCoT + OC-MERO + CRISP. \\
\bottomrule
\end{tabular}
\end{table*}

\subsection{Metrics}
\label{subsec:metrics}

Let $a_i$ be the selected action for root $i$ and $a_{i,\mathrm{nom}}$ be the nominal utility maximizer.  Unless otherwise noted, metrics are averaged over root scenes with bootstrap confidence intervals.

\paragraph{Observation-consistent recovery success.}
Expected class-consistent recovery success is
\begin{equation}
\label{eq:metric_ocrs}
    \mathrm{OCRS}=\frac{1}{N}\sum_i\sum_{m=1}^{M}\hat p_i^m
    Y_{i,a_i,\mathrm{oc}}^{m,\star}.
\end{equation}
It measures whether the selected prefix preserves a witness that can be chosen from post-prefix observations.

\paragraph{Selected lower-tail recoverability and false admission.}
\begin{align}
\label{eq:metric_slr}
    \mathrm{SLR}&=\frac{1}{N}\sum_i R_{i,a_i}^{\star},\\
\label{eq:metric_fr}
    \mathrm{FR}&=\frac{1}{N}\sum_i \mathbf 1[R_{i,a_i}^{\star}<\eta_R].
\end{align}
FR is the primary selected-action calibration metric: it counts prefixes admitted by the planner whose observation-consistent teacher recoverability is below the target threshold.

\paragraph{Oracle leakage gap.}
\begin{equation}
\label{eq:metric_olg}
    \mathrm{OLG}=\frac{1}{N}\sum_i
    \left(R_{i,a_i}^{\mathrm{oracle}}-R_{i,a_i}^{\star}\right).
\end{equation}
A large OLG indicates that oracle existential labels overstate deployable recovery by changing witnesses across indistinguishable futures.

\paragraph{Same-root ranking and regret.}
\begin{align}
\label{eq:metric_pra}
    \mathrm{PRA}&=
    \frac{\sum_i\sum_{a<b}\mathbf 1[\mathrm{sign}(\hat R_{i,a}-\hat R_{i,b})=\mathrm{sign}(R_{i,a}^{\star}-R_{i,b}^{\star})]}
    {\sum_i {K_i\choose2}},\\
\label{eq:metric_srr}
    \mathrm{SRR}&=\frac{1}{N}\sum_i
    \left[\max_{a\in\mathcal A_0(b_i)}R_{i,a}^{\star}-R_{i,a_i}^{\star}\right].
\end{align}
PRA tests whether root sharing improves relative action ranking; SRR measures the recoverability lost by the chosen action.

\paragraph{Evidence and profile errors.}
For signed evidence vector $\bm g$ with scale vector $\bm\sigma_g$,
\begin{equation}
\label{eq:metric_gme}
    \mathrm{GME}=\frac{1}{|\Omega_G|}\sum_{(i,a,j,m)\in\Omega_G}
    \left\|\frac{\widehat{\bm g}_{i,a,j}^{m}-\bm g_{i,a,j}^{m,\star}}{\bm\sigma_g}\right\|_1.
\end{equation}
The recoverability profile error is
\begin{equation}
\label{eq:metric_rpe}
    \mathrm{RPE}=\frac{1}{|\Omega_A|}\sum_{(i,a)\in\Omega_A}|\hat R_{i,a}-R_{i,a}^{\star}|.
\end{equation}

\paragraph{Witness and observation consistency.}
Witness accuracy is evaluated only when the best class witness exceeds the runner-up by at least $\epsilon_w$:
\begin{equation}
\label{eq:metric_wa}
    \mathrm{WA}=\frac{1}{|\Omega_W|}\sum_{(i,a,m)\in\Omega_W}
    \mathbf 1[\hat j_{i,a}^{m}=j_{i,a,c_{i,a}(m)}^{\star}].
\end{equation}
Observation-consistency violation is
\begin{equation}
\label{eq:metric_ocv}
    \mathrm{OCV}=\frac{1}{|\Omega_O|}\sum_{(i,a,m,m')\in\Omega_O}
    D_{\mathrm{JS}}(\mu_{i,a}^{m}\Vert\mu_{i,a}^{m'}),
\end{equation}
where $\Omega_O$ contains equivalent post-prefix observation pairs.

\paragraph{Closed-loop safety, utility, and post-contact metrics.}
We report collision rate (CR), at-fault collision rate (AFCR), off-road rate (OR), wrong-way rate (WW), time-to-collision violation (TTC-V), comfort violation (CV), and route-progress ratio
\begin{equation}
\label{eq:metric_progress}
    \mathrm{Prog}=\frac{1}{N}\sum_i
    \frac{s_i^{\mathrm{final}}-s_i^0}{s_i^{\mathrm{goal}}-s_i^0+\epsilon}.
\end{equation}
For contact roots, post-contact stabilization success is
\begin{equation}
\label{eq:metric_pcss}
    \mathrm{PCSS}=\frac{1}{N_{\mathrm{contact}}}\sum_{i\in\mathrm{contact}}
    \mathbf 1[\neg\mathrm{secondary}_i\wedge\mathrm{stable}_i\wedge(\mathrm{safe\_stop}_i\vee\mathrm{rejoin}_i)].
\end{equation}

\paragraph{Harm non-inferiority and minimal intervention.}
\begin{align}
\label{eq:metric_hniv}
    \mathrm{HNIV}&=\frac{1}{N_{\mathrm{contact}}}\sum_{i\in\mathrm{contact}}
    \mathbf 1[H_{i,a_i}^{\star}>\min_{a\in\mathcal A_0(b_i)}H_{i,a}^{\star}+\epsilon_H],\\
\label{eq:metric_mir}
    \mathrm{MIR}&=\frac{1}{N_{\mathrm{nomrec}}}\sum_{i\in\mathrm{nomrec}}
    [U_{\mathrm{drv}}(b_i,a_{i,\mathrm{nom}})-U_{\mathrm{drv}}(b_i,a_i)]_+ .
\end{align}
HNIV detects whether recovery is bought by increasing immediate harm; MIR measures utility loss on scenes where the nominal action was already recoverable.

\subsection{Closed-Loop Baseline Comparison}
\label{subsec:main_results}

\begin{table*}[t]
\centering
\footnotesize
\caption{Main closed-loop comparison on Test-ID.  OCRS, SLR, PRA, PCSS, and Prog should be higher; FR, OLG, SRR, HNIV, CR, OR, and latency should be lower.}
\label{tab:main_closed_loop}
\begin{tabular}{lcccccccccccc}
\toprule
Method & OCRS $\uparrow$ & FR $\downarrow$ & SLR $\uparrow$ & OLG $\downarrow$ & PRA $\uparrow$ & SRR $\downarrow$ & PCSS $\uparrow$ & HNIV $\downarrow$ & CR $\downarrow$ & OR $\downarrow$ & Prog $\uparrow$ & ms $\downarrow$ \\
\midrule
NOM & TBD & TBD & TBD & TBD & -- & TBD & TBD & TBD & TBD & TBD & TBD & TBD \\
Risk-CVaR & TBD & TBD & TBD & TBD & TBD & TBD & TBD & TBD & TBD & TBD & TBD & TBD \\
PSF-Backup & TBD & TBD & TBD & TBD & TBD & TBD & TBD & TBD & TBD & TBD & TBD & TBD \\
Contingency & TBD & TBD & TBD & TBD & TBD & TBD & TBD & TBD & TBD & TBD & TBD & TBD \\
WM-Recovery & TBD & TBD & TBD & TBD & TBD & TBD & TBD & TBD & TBD & TBD & TBD & TBD \\
Scalar-R & TBD & TBD & TBD & TBD & TBD & TBD & TBD & TBD & TBD & TBD & TBD & TBD \\
Viable-Hand & TBD & TBD & TBD & TBD & TBD & TBD & TBD & TBD & TBD & TBD & TBD & TBD \\
Oracle-OC & TBD & TBD & TBD & 0.000 & 1.000 & 0.000 & TBD & TBD & TBD & TBD & TBD & -- \\
Ours & TBD & TBD & TBD & TBD & TBD & TBD & TBD & TBD & TBD & TBD & TBD & TBD \\
\bottomrule
\end{tabular}
\end{table*}

\begin{table*}[t]
\centering
\footnotesize
\caption{Closed-loop results by root tag.  The same OC-RAP policy and CRISP selector are used for every tag.}
\label{tab:regime_results}
\begin{tabular}{llcccccccc}
\toprule
Tag & Method & OCRS $\uparrow$ & FR $\downarrow$ & SLR $\uparrow$ & CR $\downarrow$ & AFCR $\downarrow$ & PCSS $\uparrow$ & HNIV $\downarrow$ & Prog $\uparrow$ \\
\midrule
Normal & NOM & TBD & TBD & TBD & TBD & TBD & -- & -- & TBD \\
Normal & Ours & TBD & TBD & TBD & TBD & TBD & -- & -- & TBD \\
Low-headroom & Risk-CVaR & TBD & TBD & TBD & TBD & TBD & -- & -- & TBD \\
Low-headroom & Ours & TBD & TBD & TBD & TBD & TBD & -- & -- & TBD \\
Near-contact & PSF-Backup & TBD & TBD & TBD & TBD & TBD & -- & TBD & TBD \\
Near-contact & Ours & TBD & TBD & TBD & TBD & TBD & -- & TBD & TBD \\
Contact/post-contact & Contingency & TBD & TBD & TBD & TBD & TBD & TBD & TBD & TBD \\
Contact/post-contact & Ours & TBD & TBD & TBD & TBD & TBD & TBD & TBD & TBD \\
\bottomrule
\end{tabular}
\end{table*}

\subsection{Offline Recovery Inference}
\label{subsec:offline_results}

\begin{table}[t]
\centering
\caption{Offline recovery inference, witness prediction, and observation consistency.}
\label{tab:offline_recovery}
\begin{tabular}{lccccc}
\toprule
Method & GME $\downarrow$ & RPE $\downarrow$ & PRA $\uparrow$ & WA $\uparrow$ & OCV $\downarrow$ \\
\midrule
WM-Recovery & -- & TBD & TBD & TBD & TBD \\
Scalar-R & -- & TBD & TBD & -- & -- \\
Viable-Hand & TBD & TBD & TBD & TBD & -- \\
ReCoT only & TBD & TBD & TBD & TBD & TBD \\
ReCoT + black-box agg. & TBD & TBD & TBD & TBD & TBD \\
ReCoT + OC-MERO & TBD & TBD & TBD & TBD & TBD \\
\bottomrule
\end{tabular}
\end{table}

\subsection{Ablations}
\label{subsec:ablations}

\begin{table*}[t]
\centering
\footnotesize
\caption{Component ablations.  Rows isolate root sharing, observation consistency, $\mu$-weighted existential pooling, monotonicity, lower-tail aggregation, and calibrated selection.}
\label{tab:component_ablation}
\begin{tabular}{lcccccccc}
\toprule
Variant & OCRS $\uparrow$ & FR $\downarrow$ & SLR $\uparrow$ & OLG $\downarrow$ & PRA $\uparrow$ & OCV $\downarrow$ & HNIV $\downarrow$ & MIR $\downarrow$ \\
\midrule
Full OC-RAP & TBD & TBD & TBD & TBD & TBD & TBD & TBD & TBD \\
No root sharing & TBD & TBD & TBD & TBD & TBD & TBD & TBD & TBD \\
Oracle existential labels & TBD & TBD & TBD & TBD & TBD & TBD & TBD & TBD \\
No observation-consistency loss & TBD & TBD & TBD & TBD & TBD & TBD & TBD & TBD \\
Mean-over-options pooling & TBD & TBD & TBD & TBD & TBD & TBD & TBD & TBD \\
Unweighted normalized logsumexp & TBD & TBD & TBD & TBD & TBD & TBD & TBD & TBD \\
Unconstrained MLP calibrator & TBD & TBD & TBD & TBD & TBD & TBD & TBD & TBD \\
Expectation over modes & TBD & TBD & TBD & TBD & TBD & TBD & TBD & TBD \\
No CRISP calibration & TBD & TBD & TBD & TBD & TBD & TBD & TBD & TBD \\
No harm constraints & TBD & TBD & TBD & TBD & TBD & TBD & TBD & TBD \\
No rule/hard-shell constraint & TBD & TBD & TBD & TBD & TBD & TBD & TBD & TBD \\
\bottomrule
\end{tabular}
\end{table*}

\begin{table}[t]
\centering
\caption{Input and resolution ablations.}
\label{tab:bev_ablation}
\begin{tabular}{lcccc}
\toprule
Input & GME $\downarrow$ & PRA $\uparrow$ & FR $\downarrow$ & ms $\downarrow$ \\
\midrule
Vector-only & TBD & TBD & TBD & TBD \\
BEV-only, $128^2$ & TBD & TBD & TBD & TBD \\
BEV-only, $200^2$ & TBD & TBD & TBD & TBD \\
BEV-only, $256^2$ & TBD & TBD & TBD & TBD \\
BEV-only, $320^2$ & TBD & TBD & TBD & TBD \\
BEV + ego history & TBD & TBD & TBD & TBD \\
BEV + ego/agent history & TBD & TBD & TBD & TBD \\
Full BEV + uncertainty & TBD & TBD & TBD & TBD \\
\bottomrule
\end{tabular}
\end{table}

\subsection{Calibration and Reliability}
\label{subsec:calibration_results}

\begin{table}[t]
\centering
\caption{Selected-action calibration.  Target risks are $\delta_R,\delta_H,\delta_\Delta,\delta_C$.}
\label{tab:calibration_results}
\begin{tabular}{lccccc}
\toprule
Selector & $q_R$ & $q_H$ & $q_\Delta$ & $q_C$ & Test FR \\
\midrule
Uncalibrated & $0$ & $0$ & $0$ & $0$ & TBD \\
Score calibration only & TBD & TBD & TBD & TBD & TBD \\
Selected-action CRC & TBD & TBD & TBD & TBD & TBD \\
Mondrian selected-action CRC & TBD & TBD & TBD & TBD & TBD \\
\bottomrule
\end{tabular}
\end{table}

\begin{table}[t]
\centering
\caption{Constraint-specific selected-action false admission rates.}
\label{tab:constraint_calibration}
\begin{tabular}{lcccc}
\toprule
Method & $R$ violation & $H$ violation & $\Delta H$ violation & $C$ violation \\
\midrule
Uncalibrated & TBD & TBD & TBD & TBD \\
Selected-action CRC & TBD & TBD & TBD & TBD \\
Mondrian CRC & TBD & TBD & TBD & TBD \\
\bottomrule
\end{tabular}
\end{table}

\subsection{Robustness, Domain Shift, and Runtime}
\label{subsec:robustness_runtime}

We evaluate robustness under low friction, actuation delay, localization noise, BEV channel dropout, occlusion release, map staleness, traffic-density shift, and surrounding-agent aggressiveness.  Test-Shift changes these factors outside the training range while preserving the same label definitions.

\begin{table*}[t]
\centering
\footnotesize
\caption{Robustness and domain-shift results.}
\label{tab:robustness}
\begin{tabular}{llcccccc}
\toprule
Shift & Method & OCRS $\uparrow$ & FR $\downarrow$ & CR $\downarrow$ & PCSS $\uparrow$ & HNIV $\downarrow$ & Prog $\uparrow$ \\
\midrule
Low friction & Risk-CVaR & TBD & TBD & TBD & TBD & TBD & TBD \\
Low friction & Ours & TBD & TBD & TBD & TBD & TBD & TBD \\
Actuation delay & PSF-Backup & TBD & TBD & TBD & TBD & TBD & TBD \\
Actuation delay & Ours & TBD & TBD & TBD & TBD & TBD & TBD \\
Occlusion release & WM-Recovery & TBD & TBD & TBD & TBD & TBD & TBD \\
Occlusion release & Ours & TBD & TBD & TBD & TBD & TBD & TBD \\
High traffic density & Contingency & TBD & TBD & TBD & TBD & TBD & TBD \\
High traffic density & Ours & TBD & TBD & TBD & TBD & TBD & TBD \\
Map staleness & Scalar-R & TBD & TBD & TBD & TBD & TBD & TBD \\
Map staleness & Ours & TBD & TBD & TBD & TBD & TBD & TBD \\
\bottomrule
\end{tabular}
\end{table*}

\begin{table}[t]
\centering
\caption{Runtime decomposition in milliseconds per planning step.}
\label{tab:runtime}
\begin{tabular}{lccccc}
\toprule
Setting & BEV & ReCoT & OC-MERO & CRISP & Total \\
\midrule
$K=16,M=4,L\le6$ & TBD & TBD & TBD & TBD & TBD \\
$K=32,M=8,L\le12$ & TBD & TBD & TBD & TBD & TBD \\
$K=64,M=8,L\le12$ & TBD & TBD & TBD & TBD & TBD \\
$K=32$ with refinement & TBD & TBD & TBD & TBD & TBD \\
\bottomrule
\end{tabular}
\end{table}

\appendix
\section{Implementation and Reproducibility Details}
\label{app:implementation_new}

This appendix specifies the objects used by the method: ego dynamics and prefix projection, recovery affordance tokens, teacher rollouts and observation-consistent labels, ReCoT implementation, calibrated selection, and the MetaDrive--Waymo recovery dataset.  All thresholds are in SI units unless noted otherwise.

\subsection{Dynamics, Prefix Parameterization, and Projection}
\label{app:action_parameterization_new}

The online planner uses a kinematic bicycle model with sampling interval $\Delta t$:
\begin{equation}
\label{eq:app_bicycle_new}
\begin{split}
    x_{t+1}&=x_t+\Delta t\,v_t\cos(\psi_t+\beta_t),\\
    y_{t+1}&=y_t+\Delta t\,v_t\sin(\psi_t+\beta_t),\\
    \psi_{t+1}&=\psi_t+\Delta t\,{v_t}\sin\beta_t/l_r,\\
    v_{t+1}&=\mathrm{clip}(v_t+\Delta t\,a_t,0,v_{\max}),
\end{split}
\end{equation}
where $\beta_t=\arctan(l_r\tan\delta_t/(l_f+l_r))$.  Teacher validation in high-slip and post-contact states additionally rolls out a dynamic single-track model with tire saturation and impulse perturbations.  The admissible input set is
\begin{equation}
\label{eq:app_input_constraints_new}
    |a_t|\le a_{\max},\quad |j_t|\le j_{\max},\quad
    |\kappa_t|\le\kappa_{\max},\quad |\dot\kappa_t|\le\dot\kappa_{\max}.
\end{equation}

A prefix parameter $\xi$ contains terminal Frenet offset, terminal speed, acceleration knots, and curvature knots over $H_p$ steps.  The projection $a=\Gamma(\xi)$ solves a quadratic program that minimizes deviation from the neural proposal while enforcing Eq.~\eqref{eq:app_bicycle_new}, Eq.~\eqref{eq:app_input_constraints_new}, static swept-geometry clearance, route-compatible heading, and initial-state continuity.  Prefixes that remain infeasible after projection are discarded before recovery evaluation.

\subsection{Recovery Affordance Token Specification}
\label{app:recovery_affordances}

Each recovery affordance token has the fields in Eq.~\eqref{eq:method_affordance_set}.  The implementation uses the following tag set only for interpretability and data balance:
\begin{equation}
\label{eq:app_tag_set_new}
    d_j\in\{\mathrm{lane},\mathrm{stop},\mathrm{yield},\mathrm{gap},\mathrm{escape},\mathrm{rejoin},\mathrm{stabilize}\}.
\end{equation}
The online algorithm never branches on this tag; all tokens are scored by the same evidence head and monotone operator.

\paragraph{Anchor generation.}
For each prefix, anchors are generated from the route lane graph, legal stop regions, protected conflict-zone boundaries, visible free-space pockets, route-compatible rejoin lanes, and post-contact stabilization basins.  An anchor is retained if a relaxed kinodynamic search can connect the post-prefix ego state to the anchor within $H_r$ steps while respecting coarse drivable-space and input bounds.  At most $L$ tokens with largest relaxed reachability margin and spatial diversity are kept.

\paragraph{Hard shell.}
For state $x$, belief $b$, and token $w_j$, the physical shell is
\begin{equation}
\label{eq:app_phys_shell_new}
\begin{split}
    h_{\mathrm{phys}}(x,b,w_j)=\min\{&d_{\mathrm{coll}},d_{\mathrm{road}},m_{\mathrm{ctrl}},m_{\mathrm{fric}},\\
    &m_{\mathrm{yaw}},m_{\mathrm{side}},m_{\mathrm{secondary}}\},
\end{split}
\end{equation}
where $d_{\mathrm{coll}}$ is signed swept-footprint clearance to occupied or protected space; $d_{\mathrm{road}}$ is signed distance to drivable boundary; $m_{\mathrm{ctrl}}$ is distance to input limits; $m_{\mathrm{fric}}$ is friction-circle margin; $m_{\mathrm{yaw}}$ and $m_{\mathrm{side}}$ bound yaw-rate and sideslip; and $m_{\mathrm{secondary}}$ measures clearance to secondary-collision hazards after first contact.  The rule shell is
\begin{equation}
\label{eq:app_rule_shell_new}
    h_{\mathrm{rule}}(x,b,w_j)=\min\{m_{\mathrm{lane}},m_{\mathrm{signal}},m_{\mathrm{stop}},m_{\mathrm{yield}},m_{\mathrm{protected}}\}.
\end{equation}
Rule slack is zero for no-contact viable recovery.  It is allowed only in minimum-risk fallback and is recorded as $C_a$ for calibration.

\paragraph{Conservative recovery potential.}
Each token has
\begin{equation}
\label{eq:app_recovery_potential_new}
    \overline V_j(x,b)=V_j^{\mathrm{ana}}(x,b)+\Delta V_{\psi_j}(x,b)+q_{V,j},
\end{equation}
where $V_j^{\mathrm{ana}}$ is an analytic anchor distance or stabilize-avoid potential, $\Delta V_{\psi_j}$ is the learned residual, and $q_{V,j}$ is a calibration residual chosen on held-out rollouts.  A state is inside the token recovery set when
\begin{equation}
\label{eq:app_token_set_new}
    x\in\mathcal C_j(b)\iff
    h_{\mathrm{phys}}(x,b,w_j)\ge0,
    \quad h_{\mathrm{rule}}(x,b,w_j)+s^{\mathrm{rule}}\ge0,
    \quad \overline V_j(x,b)\le0.
\end{equation}
The attached fallback controller must also decrease the potential over one step:
\begin{equation}
\label{eq:app_decrease_new}
    \Delta_{\kappa_j}\overline V_j(x,b)=
    \overline V_j(f(x,\kappa_j(x,b)),b^+)-\overline V_j(x,b)\le-\rho_j.
\end{equation}

\paragraph{Hybrid recovery specifications.}
The same token set supports three specifications.  No-contact viable recovery is
\begin{equation}
\label{eq:app_g_no_new}
    G_{\mathrm{no}}=\min_{\tau\le H_r}\{h_{\mathrm{phys}},h_{\mathrm{rule}},-\overline V_j, -\Delta_{\kappa_j}\overline V_j-\rho_j,h_{\mathrm{hold}}\}_{\tau}.
\end{equation}
Minimum-risk contact-imminent recovery is
\begin{equation}
\label{eq:app_g_mr_new}
    G_{\mathrm{mr}}=\min\{g_{\mathrm{harm\text{-}noninf}},g_{\mathrm{secondary}},g_{\mathrm{ctrl}},g_{\mathrm{road}},g_{\mathrm{fallback}}\}.
\end{equation}
Post-contact stabilization is
\begin{equation}
\label{eq:app_g_post_new}
    G_{\mathrm{post}}=\min\{g_{\mathrm{stable}},g_{\mathrm{secondary}},g_{\mathrm{safe\text{-}stop}},g_{\mathrm{road}},g_{\mathrm{ctrl}}\}.
\end{equation}
The reported teacher margin $G_{a,j}^{m,\star}$ in Eq.~\eqref{eq:method_margin} is the largest specification margin available to the witness under the simulator state sequence, with rule slack accounted for only in the minimum-risk specification.  Equivalently, Eq.~\eqref{eq:method_margin} uses a maximum over recovery specifications, while each individual specification remains a minimum over its own hard-shell and return conditions.  The specification identity is stored as supervision but is not used as an online regime switch.

\subsection{Teacher Rollouts, Root Sharing, and Observation-Consistent Labels}
\label{app:teacher_labels_new}

For each root scene $i$, sample $M$ root seeds $\omega_i^m$.  A seed contains other-agent intent class, reaction-time perturbation, aggressiveness, occlusion-release time, friction, actuation delay, perception noise, and simulator stochasticity.  The same $\omega_i^m$ is reused for every candidate prefix in the root scene.  A teacher rollout for $(a,j,m)$ executes prefix $a$ for $H_p$ steps, switches to the fallback controller $\kappa_j$, and evaluates the signed margins in Appendix~\ref{app:recovery_affordances} over $H_r$ recovery steps.  The option label is
\begin{equation}
\label{eq:app_option_label_new}
    Y_{a,j}^{m,\star}=\mathbf 1[G_{a,j}^{m,\star}\ge0].
\end{equation}
Harm $H_a^{m,\star}$ is computed at first contact or closest unavoidable approach, using relative speed, overlap proxy, vulnerable-road-user weight, and secondary-collision exposure.  The post-contact deficit $K_{a,j}^{m,\star}$ is the positive part of stabilization, road-boundary, and secondary-collision margins after contact.

Executing prefix $a$ under root $m$ yields post-prefix observation $o_a^m$.  The teacher posterior label is
\begin{equation}
\label{eq:app_posterior_label_new}
    \beta_{a}^{m,\star}(\ell)=
    \frac{\hat p^\ell\mathbf 1[d_{\mathrm{obs}}(o_a^m,o_a^\ell)\le\epsilon_o]}
    {\sum_{r=1}^{M}\hat p^r\mathbf 1[d_{\mathrm{obs}}(o_a^m,o_a^r)\le\epsilon_o]}.
\end{equation}
The distance $d_{\mathrm{obs}}$ is a weighted sum of visible-agent pose distance, occupancy IoU distance, traffic-control disagreement, occlusion mask disagreement, and ego-state difference.  Modes with mutual distance below $\epsilon_o$ form an observation class $c_a(m)$.

The observation-consistent witness label is computed per observation class:
\begin{equation}
\label{eq:app_class_witness_new}
    j_c^\star=\argmax_j \sum_{m:c_a(m)=c}\hat p^mY_{a,j}^{m,\star},
\end{equation}
with ties broken by larger mean margin, smaller harm, smaller post-contact deficit, and finally smaller controller effort.  The observation-consistent mode label is
\begin{equation}
\label{eq:app_oc_label_new}
    Y_{a,\mathrm{oc}}^{m,\star}=Y_{a,j_{c_a(m)}^\star}^{m,\star}.
\end{equation}
The observation-consistent teacher recoverability used in experiments is
\begin{equation}
\label{eq:app_true_recoverability_new}
    R_a^\star=\mathrm{LCVaR}_{\alpha_R}\left(\{Y_{a,\mathrm{oc}}^{m,\star},\hat p^m\}_{m=1}^{M}\right).
\end{equation}
For binary labels, maximizing Eq.~\eqref{eq:app_class_witness_new} maximizes the lower-tail objective because lower-tail binary recoverability is a monotone function of weighted success probability.  Same-root ranking pairs are formed only within the same root scene.

\subsection{ReCoT Architecture and Training Losses}
\label{app:recot_architecture_new}

The scene encoder has two branches.  The vector branch encodes lane polylines, road boundaries, traffic-control elements, crosswalks, and actor tracks with a polyline transformer.  The BEV branch encodes occupancy, drivable area, route, occlusion, uncertainty, and ego history with a four-stage residual CNN.  The fused scene tokens are cross-attended by prefix, root, and recovery-affordance tokens.

Prefix tokens are sampled from $H_p+1$ ego poses, controls, curvature, jerk, swept boxes, terminal Frenet state, and route-progress features.  Root tokens contain learned seed embeddings plus perturbation metadata used by the teacher.  Recovery-affordance tokens contain anchor polyline features, controller parameters, hard-shell slack at the post-prefix state, relaxed reachability margin, and anchor type embedding.  The tuple representation in Eq.~\eqref{eq:method_recot_tuple} is produced by two attention blocks alternating over the $(a,m)$ plane and the $(a,j)$ plane, followed by a small MLP head for each evidence component.

The Huber margin loss is applied componentwise to the signed vector
\begin{equation}
\label{eq:app_margin_vector_new}
    \bm g_{a,j}^{m,\star}=(G_{\mathrm{no}},G_{\mathrm{mr}},G_{\mathrm{post}},h_{\mathrm{phys}},h_{\mathrm{rule}},-\overline V_j,h_{\mathrm{hold}},g_{\mathrm{secondary}})_{a,j}^{m,\star}.
\end{equation}
Harm and post-contact deficit heads use pinball losses at quantiles $\{0.5,0.8,0.9\}$.  The uncertainty head predicts a Gaussian or Student-$t$ residual over evidence errors and is trained by negative log-likelihood.  The monotone calibrator in Eq.~\eqref{eq:method_option_score_new} is implemented as a nonnegative-weight MLP or a deep lattice network~\cite{you2017deep}.  Its sampled monotonicity loss is
\begin{equation}
\label{eq:app_mono_loss_new}
    \mathcal L_{\mathrm{mono}}=
    \sum_{(x,x^+)}[f_\phi^{\mathrm{mono}}(x)-f_\phi^{\mathrm{mono}}(x^+)+\gamma_m]_+,
\end{equation}
where $x^+$ improves one safety margin or reduces one deficit while holding other inputs fixed.  Same-root ranking uses
\begin{equation}
\label{eq:app_rank_loss_new}
    \ell_{\mathrm{rank}}=
    \sum_{a_i,a_j}\mathbf 1[R_{a_i}^\star>R_{a_j}^\star]
    [\gamma_R-(R_{a_i}-R_{a_j})]_+.
\end{equation}

\subsection{Calibration and Online Selection Algorithm}
\label{app:calibration_selector_new}

Calibration uses held-out root scenes disjoint from training and validation roots.  The target risk is selected-action false recovery admission: a selected prefix is falsely admitted when the calibrated set accepts it but the observation-consistent teacher label is below threshold.  For recovery, define the nonconformity score
\begin{equation}
\label{eq:app_nonconf_recovery_new}
    s_i^R=\eta_R-R_{a_i}\quad\text{on examples with }R_{a_i}^\star<\eta_R.
\end{equation}
The offset $q_R$ is the smallest value satisfying the empirical conformal risk-control constraint at level $\delta_R$.  Harm, harm-gap, and rule offsets are computed analogously from $H_a-H_a^\star$, $\Delta H_a-\Delta H_a^\star$, and $C_a-C_a^\star$.  Mondrian calibration strata are source domain, recovery tag, and root criticality; if a stratum has fewer than $n_{\min}$ examples, the global offset is used.

\begin{algorithm}[t]
\caption{OC-RAP online selection}
\label{alg:ocrap_selector}
\begin{algorithmic}[1]
\Require belief $b_t$, route $q_t$, calibration offsets $q$
\State Generate prefix proposals $\mathcal A_0(b_t)$ and project with $\Gamma$
\State Optionally refine proposals by Eq.~\eqref{eq:method_refinement_new}; retain $\mathcal A$
\For{$a\in\mathcal A$}
    \State Generate recovery-affordance tokens $\mathcal W_a(b_t)$
    \State Predict root-shared evidence $\widehat e_{a,j}^{m}$ for all retained $(j,m)$
    \State Compute OC-MERO profile $\Pi_a$ by Eqs.~\eqref{eq:method_existential_new}--\eqref{eq:method_profile_new}
\EndFor
\State $a_{\mathrm{nom}}\gets\argmax_{a\in\mathcal A}U_{\mathrm{drv}}(b_t,a)$
\State Construct $\mathcal S(q)$ by Eq.~\eqref{eq:method_admissible_set}
\If{$a_{\mathrm{nom}}\in\mathcal S(q)$}
    \State $a^\star\gets a_{\mathrm{nom}}$; $s\gets\mathrm{nominal\ accepted}$
\ElsIf{$\mathcal S(q)\ne\emptyset$}
    \State $a^\star\gets$ Eq.~\eqref{eq:method_constrained_selector}; $s\gets\mathrm{calibrated\ switch}$
\Else
    \State $a^\star\gets$ Eq.~\eqref{eq:method_relaxation}; $s\gets\mathrm{minimum\ risk}$
\EndIf
\State Attach $\kappa_{\hat w_{a^\star}}$ if $\hat w_{a^\star}$ exists; otherwise attach stop/stabilize fallback
\State \Return $(a^\star,\Pi_{a^\star},\hat w_{a^\star},\kappa,s)$
\end{algorithmic}
\end{algorithm}

\subsection{MetaDrive--Waymo Recovery Dataset Construction}
\label{app:dataset_construction_new}

The dataset combines procedural MetaDrive scenes with Waymo Open Motion Dataset (WOMD) scenes converted to an interactive MetaDrive-compatible format.  MetaDrive is compositional and supports procedural scenario generation and real-data replay~\cite{li2022metadrive}.  WOMD provides object tracks and map data, with segments recorded at $10\,\mathrm{Hz}$ and scenario windows containing history and future trajectories~\cite{ettinger2021large}.  ScenarioNet-style conversion is used to normalize heterogeneous traffic scenarios into a unified object-centric format that can be replayed and interacted with in MetaDrive~\cite{li2023scenarionet}.

\paragraph{Sources.}
We use three source domains.
\begin{enumerate}[leftmargin=*]
    \item \textbf{Procedural MetaDrive.} Maps are generated from randomized road blocks, lane counts, intersections, merges, roundabouts, traffic controls, and traffic densities.  Root criticality is controlled by traffic density, time-to-conflict, occlusion release, and drivable-space clearance.
    \item \textbf{WOMD replay.} Scenario protocol buffers are parsed into lane polylines, road edges, crosswalks, stop signs, traffic signals, object tracks, object types, validity masks, and the self-driving-car track when available.  Tracks are downsampled from $10\,\mathrm{Hz}$ to the planning rate.  Other agents are replayed by log-tracking controllers with perturbable reaction and intent parameters, allowing closed-loop ego interaction without changing the original map geometry.
    \item \textbf{Hybrid stress.} WOMD maps and traffic initializations are used as seeds, then conflict timing, occluded actor release, lead braking, cut-in aggressiveness, friction, and actuation delay are perturbed to increase low-headroom, near-contact, and post-contact coverage.
\end{enumerate}

\paragraph{Root extraction.}
A root scene is defined by a scenario identifier, ego state at planning time, local map crop, route command, visible history, hidden future root seeds, and simulator random seed.  For WOMD, the ego is the SDC track when valid; otherwise, an interactive target track with sufficient map support is selected and treated as ego for offline simulation.  Roots are rejected if the ego is off-map, if fewer than three relevant dynamic agents are available within $60\,\mathrm{m}$ in interactive tags, if map lanes cannot be connected to a route corridor, or if all candidate prefixes are trivially infeasible.

\paragraph{History and observations.}
The canonical input uses $H_h=5$ frames at $\Delta t=0.2\,\mathrm{s}$, matching one second of history after WOMD downsampling.  Procedural scenes store the same history length.  The BEV crop is $80\,\mathrm{m}\times80\,\mathrm{m}$ at $256\times256$ resolution.  Channels include drivable area, lane centerlines, lane boundaries, route corridor, traffic control, static obstacles, visible dynamic occupancy, dynamic-agent velocity, ego footprint history, current ego footprint, occlusion mask, uncertainty, free-space-pocket anchors, and recovery-affordance anchors.  Structured tensors additionally store vectorized map elements, actor histories, prefix control points, swept polygons, and token fields.

\paragraph{Criticality tags.}
Each root is assigned one balancing tag but no online gate uses this tag:
\begin{equation}
\label{eq:app_tags_new}
    \mathrm{tag}\in\{\mathrm{normal},\mathrm{low\text{-}headroom},\mathrm{near\text{-}contact},\mathrm{post\text{-}contact}\}.
\end{equation}
Normal roots have high time-to-conflict and at least one high-margin route-following witness.  Low-headroom roots are collision-free but have limited braking, lateral escape, or route-rejoin margin.  Near-contact roots include cut-in, lead braking, merge, crossing, occluded release, or forced-yield conflicts with small time-to-conflict.  Post-contact roots disable crash termination and include unavoidable or already-contacted states for stabilization and secondary-collision labeling.

\paragraph{Label generation.}
For each root, generate $K_{\mathrm{raw}}$ proposals from imitation, lattice, emergency, and perturbation sources.  After projection, prune to $K$ prefixes by utility diversity and swept-corridor diversity.  For each prefix, generate at most $L$ recovery-affordance tokens.  For every retained $(a,j,m)$, roll out the teacher under the shared root seed $\omega^m$, record specification margins, harm, post-contact deficit, post-prefix observation, posterior class, and observation-consistent witness labels from Appendix~\ref{app:teacher_labels_new}.  Early termination is allowed only after a hard-shell margin falls below $-2.0\,\mathrm{m}$ or a reachability lower bound proves that the token cannot be reached within the remaining horizon.

\paragraph{Splits and leakage control.}
Splits are by root scene and by original WOMD scenario identifier.  No map seed, WOMD scenario id, root time, or latent seed is shared across train, validation, calibration, and test.  WOMD challenge test labels are not used.  The shifted test split changes friction, actuation delay, localization noise, occlusion release, traffic density, and surrounding-agent aggressiveness outside the training ranges while preserving the same label definitions.

\begin{table}[t]
\centering
\caption{Default root-scene split.  Criticality tags are balanced within each source domain.}
\label{tab:app_dataset_split_new}
\begin{tabular}{lccc}
\toprule
Split & Procedural & WOMD replay & Hybrid stress \\
\midrule
Train & 10000 & 7000 & 3000 \\
Validation & 1000 & 700 & 300 \\
Calibration & 2000 & 1400 & 600 \\
Test-ID & 2000 & 1400 & 600 \\
Test-Shift & 2000 & 1400 & 600 \\
\bottomrule
\end{tabular}
\end{table}

\paragraph{Stored schema.}
Each root is stored as a compressed record with keys
\begin{equation}
\label{eq:app_schema_new}
\begin{split}
\{&\mathrm{scene\_id},\mathrm{source},\mathrm{tag},b_t,q_t,\mathcal A,\mathcal W_a,\omega^m,
\widehat p^m,\\
& G_{a,j}^{m,\star},\bm g_{a,j}^{m,\star},Y_{a,j}^{m,\star},H_a^{m,\star},K_{a,j}^{m,\star},o_a^m,\beta_a^{m,\star},j_{c}^\star,R_a^\star\}.
\end{split}
\end{equation}
Array dimensions are padded to $(K,M,L)$ with validity masks.  All physical quantities are stored in the ego frame at root time except map polylines, which additionally store their original world-frame coordinates for replay.

\subsection{Default Hyperparameters}
\label{app:hyperparameters_new}

\begin{table*}[t]
\centering
\footnotesize
\caption{Default parameters.}
\label{tab:default_hyperparameters_new}
\begin{tabular}{p{0.18\linewidth}p{0.16\linewidth}p{0.58\linewidth}}
\toprule
Symbol & Default & Meaning \\
\midrule
$\Delta t$ & $0.2\,\mathrm{s}$ & Planning and teacher rollout step. \\
$H_h$ & $5$ & Input history, one second after WOMD downsampling. \\
$H_p$ & $10$ & Prefix horizon, $2.0\,\mathrm{s}$. \\
$H_r$ & $25$ & Recovery evaluation horizon, $5.0\,\mathrm{s}$. \\
$H_{\mathrm{hold}}$ & $8$ & Dwell horizon inside a recovery token, $1.6\,\mathrm{s}$. \\
BEV extent & $80\,\mathrm{m}\times80\,\mathrm{m}$ & Ego-centered raster crop. \\
BEV resolution & $256\times256$ & Main raster resolution, $0.3125\,\mathrm{m/pixel}$. \\
$K_{\mathrm{raw}}$ & $64$ & Raw proposals before projection and pruning. \\
$K$ & $32$ & Retained prefixes per root. \\
$M$ & $8$ & Root-shared latent futures. \\
$L$ & $\le12$ & Retained recovery-affordance tokens per prefix. \\
$N_{\mathrm{ref}}$ & $5$ & Projected refinement steps. \\
$\tau_R$ & $0.20$ & Existential softmax temperature. \\
$c_R$ & $0$ & Default for $\mu$-weighted aggregation; no extra option-count normalization. \\
$\alpha_R$ & $0.20$ & Lower-tail mass for recoverability. \\
$\alpha_H,\alpha_K$ & $0.20$ & Upper-tail masses for harm and post-contact deficit. \\
$\eta_R$ & $0.70$ & Minimum calibrated lower-tail recoverability. \\
$\eta_H$ & task-defined & Maximum calibrated absolute harm exposure. \\
$\epsilon_H$ & $0.05$--$0.10$ & Relative harm non-inferiority tolerance. \\
$\delta_R,\delta_H,\delta_\Delta,\delta_C$ & $0.05$ & Target false-admission risks. \\
$n_{\min}$ & $500$ & Minimum examples for a Mondrian stratum. \\
$a_{\max}$ & $4.0\,\mathrm{m/s^2}$ & Acceleration bound. \\
$j_{\max}$ & $6.0\,\mathrm{m/s^3}$ & Jerk bound. \\
$\kappa_{\max}$ & $0.25\,\mathrm{m^{-1}}$ & Curvature bound. \\
$\dot\kappa_{\max}$ & $0.20\,\mathrm{m^{-1}s^{-1}}$ & Curvature-rate bound. \\
$\mu_{\mathrm{fric}}$ & $0.55$--$0.95$ & Training friction; shifted test extends to $0.40$--$1.05$. \\
\bottomrule
\end{tabular}
\end{table*}


\bibliographystyle{IEEEtran}
\bibliography{post-collision}

\end{document}